\chapter{Experimentação e Validação}

\section{Bases de Dados}
\label{sec:bases_de_dados}

\subsection{Bases de dados Sintéticas}

% AQUI VOU TER QUE CITAR OS AUTORES DOS FLUXOS DE DADOS
As seguintes bases de dados sintéticas foram utilizadas:
\begin{itemize}
    \item AssetNegotiationGenerator
    \item SEAGenerator
    \item RandomRBFGenerator
    \item AgrawalGenerator
    \item HyperplaneGenerator
    \item STAGGERGenerator
    \item RandomTreeGenerator
    \item WaveformGenerator
    \item LEDGenerator
\end{itemize}

\subsection{Bases de dados Reais}

% AQUI VOU TER QUE CITAR OS AUTORES DOS FLUXOS DE DADOS
A seguir, apresentamos as bases de dados reais empregadas:
\begin{itemize}
    \item AWS Spot Pricing Market
\end{itemize}

\section{Protocolo Experimental}

Utilizando as bases de dados listadas na seção \ref{sec:bases_de_dados} subtemos aos seguintes experimentos.

Inicialmente, conduzimos um experimento inserindo apenas os dados provenientes do fluxo de dados, com o objetivo de verificar se a inserção sequencial dos dados provoca o desbalanceamento da \textit{k}d-tree.

Para avaliar o desbalanceamento da árvore, consideramos o número total de nós, representado por $n$, e sua altura, denotada por $h$. A partir da quantidade de nós, é possível determinar a altura balanceada teórica da árvore, conforme a Equação~\ref{eq:altura_balanceada}.

\begin{equation}
    h_{bal} = \lfloor \log_2 n \rfloor
    \label{eq:altura_balanceada}
\end{equation}

Também comparamos essa altura com a de uma árvore rubro-negra, segundo a Equação~\ref{eq:altura_rubro_negra}.

\begin{equation}
    h_{rn} = 2 \cdot \log_2 (n + 1)
    \label{eq:altura_rubro_negra}
\end{equation}

Nos experimentos realizados, esperamos que a árvore gerada a partir de um fluxo de dados qualquer apresente altura dentro do intervalo definido pela Equação~\ref{eq:altura_ideal}.

\begin{equation}
    h_{bal} \leq h \leq h_{rn}
    \label{eq:altura_ideal}
\end{equation}

Dessa forma, é possível quantificar a qualidade estrutural das árvores geradas. Posteriormente, essa metodologia será aplicada a cenários envolvendo fluxos de dados contínuos com janelas deslizantes, onde a skd-tree será submetida.

\section{Conclusão}
